\subsection{A matching measured-information bound}
\label{sec:joint-frontier}
The quantum-relative-entropy upper bound is too large to determine the
readout gain when both the signal angle and the noise are small. A
likelihood-ratio bound combined with fidelity retains the missing angular
dependence. The ingredients are standard divergence comparison and
measurement fidelity \citep{sason2016,fuchs1995}; the following proofs
specialize them explicitly to the present model.

\begin{lemma}[Relative entropy from affinity and a ratio bound]
\label{lem:kl-affinity}
Let $p,q$ be strictly positive probability vectors with $p_y/q_y\leq R$
for every $y$, where $R\geq1$. Put $a(p,q)=\sum_y\sqrt{p_yq_y}$. Then
\begin{equation}\label{eq:kl-affinity}
 D_{\rm KL}(p\|q)\leq2(2+\log R)[1-a(p,q)].
\end{equation}
\end{lemma}
\begin{proof}
For $g(v)=v^2\log(v^2)-v^2+1$, one has $g(1)=g'(1)=0$ and
$g''(v)=4+4\log v\leq4+2\log R$ for $0<v\leq\sqrt R$.
Taylor's integral formula along the segment from $1$ to $v$ gives
$g(v)\leq(2+\log R)(v-1)^2$. Substitute
$v=\sqrt{p_y/q_y}$, multiply by $q_y$, and sum. The linear terms
cancel because both probability vectors have total mass one, while
$\sum_y(\sqrt{p_y}-\sqrt{q_y})^2=2[1-a(p,q)]$.
\end{proof}

\begin{lemma}[Fidelity and generalized eigenvalues of a phase block]
\label{lem:block-fidelity}
For a fixed depth and $\gamma>0$, abbreviate
$r=r_t$, $\vartheta=\vartheta_t$, $s=r\sin\vartheta$, and put
\begin{equation}\label{eq:eta-G}
 \eta_t=\frac{s_t^2}{1-r_t^2},\qquad
 G_t=1+\log(1+\eta_t),\qquad
 R_t=(\sqrt{1+\eta_t}+\sqrt{\eta_t})^2.
\end{equation}
Every nonzero POVM effect has outcome likelihood ratio in
$[R_t^{-1},R_t]$, and every measurement has affinity at least
\begin{equation}\label{eq:block-fidelity}
 F_t=\|\sqrt{\rho_{+,t}}\sqrt{\rho_{-,t}}\|_1
     =\sqrt{1-s_t^2}.
\end{equation}
The fidelity statement also holds at $\gamma=0$.
\end{lemma}
\begin{proof}
For any density matrices $\rho,\sigma$ and any POVM, polar decomposition
and the Hilbert--Schmidt Cauchy--Schwarz inequality give
\[
 \|\sqrt\rho\sqrt\sigma\|_1
 =\max_U|\Tr(U\sqrt\rho\sqrt\sigma)|
 \leq\sum_y|\Tr(U_\star\sqrt\rho E_y\sqrt\sigma)|
 \leq\sum_y\sqrt{\Tr(\rho E_y)\Tr(\sigma E_y)}.
\]
Here $U_\star$ is a maximizing unitary. For two-dimensional positive
matrices, the square of the trace of a positive square root is its trace
plus twice the square root of its determinant. Consequently
\[
 F_t^2=\Tr(\rho_{+,t}\rho_{-,t})
            +2\sqrt{\det\rho_{+,t}\det\rho_{-,t}}
 =\frac{1+r^2\cos(2\vartheta)}2+\frac{1-r^2}2
 =1-s^2.
\]
For full-rank states, the generalized eigenvalues have product one and
sum $2+4\eta_t$, as in Proposition~\ref{prop:residual}. They are
$R_t$ and $R_t^{-1}$. The corresponding operator inequalities
$R_t^{-1}\rho_{-,t}\preceq\rho_{+,t}\preceq R_t\rho_{-,t}$,
traced against an arbitrary positive effect, prove the ratio bound.
\end{proof}

\begin{theorem}[Uniform pointwise measured-information law]
\label{thm:pointwise-complete}
For every $\gamma>0$ and allowed depth,
\begin{equation}\label{eq:pointwise-complete}
 \frac14s_t^2G_t\leq d_M(t)
 \leq2(2+\log R_t)(1-\sqrt{1-s_t^2})
 \leq7s_t^2G_t.
\end{equation}
Thus $d_M(t)/d_t\asymp1+\log(1+\eta_t)$, with universal constants
throughout the stated model. The lower order is attained by the better
of $Y$ and the recorded random-axis readout; adaptation is unnecessary
for this constant-factor statement.
\end{theorem}
\begin{proof}
Apply Lemmas~\ref{lem:kl-affinity} and~\ref{lem:block-fidelity} to every
POVM and take the supremum. Since
$R_t\leq4(1+\eta_t)$ and $1-\sqrt{1-s_t^2}\leq s_t^2$,
\[
 2(2+\log R_t)(1-\sqrt{1-s_t^2})
 \leq2s_t^2[2+\log4+\log(1+\eta_t)]\leq7s_t^2G_t.
\]
For the lower bound, $d_t\geq2s_t^2$, and
\eqref{eq:axisinfo} gives
\[
 d_A(t)\geq\frac{s_t^2}{2r_t}
       \log\!\left(1+\frac{2(1+r_t)}{r_t}\eta_t\right)
 \geq\frac{s_t^2}{2}\log(1+\eta_t).
\]
The maximum of these two lower bounds is at least their average,
$s_t^2[1+\tfrac14\log(1+\eta_t)]\geq\tfrac14s_t^2G_t$.
Finally $2s_t^2\leq d_t\leq K_{\rm sig}s_t^2$.
\end{proof}

\subsection{An exact counting test with inconclusive observations}
\label{sec:conclusive-counting}
Stopping at the first conclusive outcome, as in
Proposition~\ref{prop:crossover}, is only one use of the three-outcome
measurement. Accumulating its signed conclusive outcomes removes that
construction's error floor at every positive noise rate.

\begin{proposition}[Repeated conclusive counting]
\label{prop:conclusive-counting}
Fix an allowed depth and $\gamma>0$. Put $q_t=\cos\vartheta_t$ and
\begin{align}
 u_t&=r_t(1-q_t)+\frac{1-r_t}{2(1+q_t)},&
 v_t&=\frac{1-r_t}{2(1+q_t)},\notag\\
 \ell_{E,t}&=\log\frac{u_t}{v_t}
 =\log\!\left(1+\frac{2r_t\sin^2\vartheta_t}{1-r_t}\right),&
 b_{E,t}&=\left\lceil\frac{\Lambda_\delta}{\ell_{E,t}}\right\rceil.
 \label{eq:conclusive-increments}
\end{align}
Apply the POVM of Proposition~\ref{prop:usd} independently to each fresh
block. Record $+1,-1,0$ for its respective outcomes $+,-,?$. Stop when
the accumulated signed count reaches either $b_{E,t}$ or $-b_{E,t}$.
The error and expected cost under either hypothesis are exactly
\begin{equation}\label{eq:conclusive-counting}
 \delta_{E,t}=\frac1{1+e^{b_{E,t}\ell_{E,t}}}\leq\delta,
 \qquad
 U_E(t,\delta)=
 \frac{c(t)b_{E,t}}{r_t(1-q_t)}
       \tanh\!\left(\frac{b_{E,t}\ell_{E,t}}2\right).
\end{equation}
Writing $U_Y(t,\delta)$ for the fixed-depth cost in
\eqref{eq:mainfrontier}, one has the explicit bound
\begin{equation}\label{eq:local-finite-upper}
 C_{M,\delta}^\star
 \leq\min_t\min\{U_Y(t,\delta),U_E(t,\delta)\}
 \leq\min_t\frac{4c(t)}{s_t^2}
              \left(1+\frac{\log(1/\delta)}{G_t}\right).
\end{equation}
\end{proposition}
\begin{proof}
Under $h=+$ the three probabilities are $u_t,v_t,1-u_t-v_t$; under
$h=-$ the first two are interchanged. Inconclusive observations have
zero log likelihood. Conditional on a conclusive observation, the signed
walk has upward probability $u_t/(u_t+v_t)$ and likelihood increment
$\ell_{E,t}$. The gambler's-ruin calculation therefore gives the stated
error. Each conclusive observation takes a geometric number of blocks
with mean $(u_t+v_t)^{-1}$. Applying the same stopping-time calculation
as in Theorem~\ref{thm:frontier} gives the displayed cost because
$u_t-v_t=r_t(1-q_t)>0$. In particular, the mean stopping time is finite.

Set $B=\log(1/\delta)$. If $\eta_t\geq1$, then
$\ell_{E,t}\geq\log(1+4\eta_t)\geq1$ and
$G_t\leq2\ell_{E,t}$. Thus $b_{E,t}\leq1+2B/G_t$ and
\[
 U_E(t,\delta)\leq\frac{2c(t)}{s_t^2}
                  \left(1+\frac{2B}{G_t}\right),
\]
where $s_t^2/[r_t(1-q_t)]=r_t(1+q_t)\leq2$.
If $\eta_t<1$, then $G_t<2$, and
\[
 U_Y(t,\delta)\leq\frac{c(t)}{2s_t^2}(B+L_{\rm sig})
 \leq\frac{c(t)}{s_t^2}\left(1+\frac{B}{G_t}\right),
\]
using $L_{\rm sig}<2$. These two bounds prove
\eqref{eq:local-finite-upper}.
\end{proof}

\subsection{The uniform noise--depth--confidence law}
\label{sec:uniform-cost}
An information-rate lower bound alone can tend to zero as noise
vanishes. A second lower bound accounts for the cost of acquiring a
decisive observation in the first place.

\begin{proposition}[An affinity converse at every noise level]
\label{prop:full-affinity}
For $0\leq\gamma\leq1$, define
\begin{equation}\label{eq:affinity-rate}
 J_F=\max_{t\in\mathcal D_\mu}
       \frac{1-\sqrt{1-s_t^2}}{c(t)},\qquad
 h_\delta=1-2\sqrt{\delta(1-\delta)}.
\end{equation}
Every admissible adaptive protocol obeys
\begin{equation}\label{eq:full-affinity-bound}
 C_{M,\delta}^\star\geq\frac{h_\delta}{J_F}.
\end{equation}
For positive noise this combines with the information converse to give
\begin{equation}\label{eq:joint-converse}
 C_{M,\delta}^\star\geq
 \max\left\{\frac{h_\delta}{J_F},\frac{k_\delta}{I_M}\right\}.
\end{equation}
\end{proposition}
\begin{proof}
For every living history, Lemma~\ref{lem:block-fidelity} gives
$a_z\geq\sqrt{1-s_{t_z}^2}$ and hence
$1-a_z\leq J_Fc(t_z)$. Substitute this inequality in the complete
stopped-affinity telescoping argument~\eqref{eq:stopped-affinity}.
Its truncation, finite-cost passage to the limit, and binary-decision
bound apply unchanged to the mixed states. Therefore
$h_\delta\leq J_F\max_h\E_h C$. The second lower bound is
Theorem~\ref{thm:measured}.
\end{proof}

\begin{theorem}[Uniform resource law for optimal single-block readout]
\label{thm:uniform-cost}
Retain the original phase cutoff, $0<\mu\leq1/8$, $D\geq1$,
$c_0\geq0$, $c_1>0$, and $0<\delta\leq1/4$. For $0<\gamma\leq1$, put
\begin{equation}\label{eq:joint-scales}
 \reff=\min\{D,\mu^{-1},\gamma^{-1}\},\quad
 A_{\rm eff}=\frac1{\mu^2}
       \left(\frac{c_0}{\reff^2}+\frac{c_1}{\reff}\right),\quad
 G_{\rm eff}=1+\log\!\left(1+\frac{\mu^2\reff}{\gamma}\right).
\end{equation}
There are universal positive constants, independent of every displayed
model parameter and of $\delta$, such that
\begin{equation}\label{eq:uniform-rates}
 J_F\asymp A_{\rm eff}^{-1},\qquad
 I_M\asymp\frac{G_{\rm eff}}{A_{\rm eff}},\qquad
 \frac{I_M}{I_Y}\asymp G_{\rm eff},
\end{equation}
and
\begin{equation}\label{eq:uniform-cost}
 \boxed{\displaystyle
 C_{M,\delta}^\star\asymp A_{\rm eff}
       \left[1+\frac{\log(1/\delta)}{G_{\rm eff}}\right].}
\end{equation}
The sharpened universal constants proved in Theorem~\ref{thm:sharp-global} give
\begin{equation}\label{eq:explicit-uniform-constants}
 \frac{A_{\rm eff}}{70}
 \left(1+\frac{\log(1/\delta)}{G_{\rm eff}}\right)
 \leq C_{M,\delta}^\star\leq
 42 A_{\rm eff}\left(1+\frac{\log(1/\delta)}{G_{\rm eff}}\right).
\end{equation}
The displayed upper bound is achieved at the single fixed depth
$t_\sharp=\min\{\overline D,\lceil\reff/2\rceil\}$, choosing $Y$ if
$\eta_{t_\sharp}<1$ and conclusive counting otherwise. For $\gamma=0$,
define $\reff=\min\{D,\mu^{-1}\}$ and $G_{\rm eff}=+\infty$,
and define the confidence-dependent term to be zero;
the cost law still holds, including $\delta=0$, by
Theorem~\ref{thm:noiseless-converse}. The finite-confidence comparison
is up to universal constants, not an exact minimax value.
\end{theorem}

The explicit upper-to-lower multiplier is 2940. These constants are
analytical and uniform, but they are not numerically sharp. Across the 246 parameter cases supplied with
Section~\ref{sec:researchresults}, the sharper converse and attainable
protocol bracket $C_{M,\delta}^\star$ between approximately $0.0584$ and
$4.312$ times the displayed scale. These are numerical evaluations on
the stated finite grid, not universal constants or interval certificates;
Figure~\ref{fig:joint-frontier} shows the brackets in a depth-limited slice.

\begin{proof}
Write $R=\reff$ temporarily and set
$t_0=\max\{1,\lfloor\pi R/8\rfloor\}$ for this scaling argument.
This depth is admissible and
$\pi R/16\leq t_0\leq R$, with $\gamma t_0\leq1$.
For every allowed $t$ with $\gamma t\leq1$, the sine bounds on
$0<2\mu t\leq\pi/4$ give
\begin{equation}\label{eq:local-effective-scales}
 s_t^2\asymp\mu^2t^2,\qquad
 \eta_t=\frac{\sin^2(2\mu t)}{e^{2\gamma t}-1}
        \asymp\frac{\mu^2t}{\gamma}.
\end{equation}
Here $e^{2u}-1\asymp u$ uniformly on $0<u\leq1$.
At $t_0$ this yields
\begin{equation}\label{eq:t0-scales}
 \frac{s_{t_0}^2}{c(t_0)}\asymp A_{\rm eff}^{-1},\qquad
 G_{t_0}\asymp G_{\rm eff}.
\end{equation}
Multiplication of the argument of $\log(1+x)$ by a fixed positive
constant changes $1+\log(1+x)$ by at most constant factors.

For the converse rate bounds, if $t\leq R$, use
$s_t^2\leq4\mu^2t^2$, $\eta_t\leq2\mu^2t/\gamma$, and the
monotonicity of $t^2/(c_0+c_1t)$ and of $\log(1+x)$ to obtain
\[
 \frac{s_t^2}{c(t)}\lesssim A_{\rm eff}^{-1},\qquad
 \frac{s_t^2G_t}{c(t)}\lesssim
                A_{\rm eff}^{-1}G_{\rm eff}.
\]
An allowed $t>R$ can occur only when $R=\gamma^{-1}$. Put
$u=\gamma t>1$. Then
\[
 \eta_t\leq\frac1{2(e^{2u}-1)},\qquad
 \frac{s_t^2}{c(t)}\leq
 \frac{4\mu^2R^2}{c(R)}u^2e^{-2u}\lesssim A_{\rm eff}^{-1}.
\]
Thus $G_t$ is uniformly bounded in this tail. Combining these estimates
with~\eqref{eq:t0-scales},
$s_t^2/2\leq1-\sqrt{1-s_t^2}\leq s_t^2$, and
Theorem~\ref{thm:pointwise-complete} proves the first two rate statements.
The last follows from $d_t\asymp s_t^2$ and the same optimization.

Uniformly for $\delta\leq1/4$,
$h_\delta\geq1-\sqrt3/2>0$ and
$k_\delta\geq\tfrac14\log(1/\delta)$. The rate bounds already proved
above and the joint converse therefore
lower-bounds the cost by a constant times
$A_{\rm eff}\max\{1,\log(1/\delta)/G_{\rm eff}\}$.
The maximum and the sum differ by at most a factor of two.
For achievability, evaluate~\eqref{eq:local-finite-upper} at $t_0$
and apply~\eqref{eq:t0-scales}. This proves the upper order uniformly
in the joint limit of noise, depth, signal, and confidence.
The improved numerical constants require keeping the depth ratio and
noise attenuation together; Theorem~\ref{thm:sharp-global} supplies that
argument and the displayed constants $1/70$ and $42$.
At zero noise
the previously proved noiseless resource law gives exactly the stated
extension.
\end{proof}

\begin{corollary}[The depth-limited regime is characterized]
\label{cor:depth-complete}
If $D\leq\min\{\mu^{-1},\gamma^{-1}\}$ and $\gamma>0$, then
\begin{align}
 \frac{I_M}{I_Y}&\asymp1+\log(1+\mu^2D/\gamma),
 \label{eq:depth-information}\\
 C_{M,\delta}^\star&\asymp
 \frac1{\mu^2}\left(\frac{c_0}{D^2}+\frac{c_1}{D}\right)
 \left[1+\frac{\log(1/\delta)}{1+\log(1+\mu^2D/\gamma)}\right].
 \label{eq:depth-complete}
\end{align}
The phase cutoff may make the largest actual depth smaller than $D$;
the constant-factor comparison remains uniform because $t_0\asymp D$
is admissible. The gain is bounded when $\mu^2D\lesssim\gamma$ and
grows logarithmically when $\mu^2D/\gamma\to\infty$.
\end{corollary}
\begin{proof}
Here $\reff=D$, so substitute into Theorem~\ref{thm:uniform-cost}.
\end{proof}

\begin{center}
\small
\begin{tabular}{p{0.26\linewidth}p{0.30\linewidth}p{0.32\linewidth}}
\toprule
Regime & Quantum ratio $I_Q/I_Y$ & Measured ratio $I_M/I_Y$\\
\midrule
Depth limited, $\reff=D$
& $1+\log(1/(\gamma D))$
& $1+\log(1+\mu^2D/\gamma)$\\
Coherent, $\reff=\mu^{-1}$
& $1+\log(\mu/\gamma)$
& $1+\log(1+\mu/\gamma)$\\
Noise limited, $\reff=\gamma^{-1}$
& $1$ & $1$\\
\bottomrule
\end{tabular}
\end{center}
All entries are up to universal constants for positive noise. The
quantum-ratio column is Proposition~\ref{prop:regimegap}; the
measured-ratio column follows from Theorem~\ref{thm:uniform-cost}.
In every row the cost is $A_{\rm eff}[1+\log(1/\delta)/G_{\rm eff}]$
up to constants. At zero noise $G_{\rm eff}=+\infty$ by convention and
the cost is comparable to $A_{\rm eff}$, including at zero error.

The additive term in~\eqref{eq:uniform-cost} is essential. It prevents
the KL-based approximation from predicting a vanishing decision cost
as the noise tends to zero. The transition occurs when
$\log(1/\delta)$ is comparable to $G_{\rm eff}$; the earlier
single-conclusive-outcome construction is a special case of the exact
counting protocol above. At fixed positive noise the leading
high-confidence constant remains exactly $1/I_M$, as in
Theorem~\ref{thm:measured}; the uniform law controls parameter limits
that this fixed-model asymptotic does not.

At positive noise, zero error with finite expected cost is impossible.
Every nonzero measurement effect has positive probability under both
hypotheses. On each event of stopping after a fixed finite number of
blocks, the transcript laws are therefore mutually absolutely continuous.
Their countable union has probability one under any finite-cost protocol,
so a terminal event cannot identify one hypothesis with certainty and
have probability zero under the other. This is consistent with the
divergence of~\eqref{eq:uniform-cost} as $\delta\downarrow 0$ at fixed
$\gamma>0$, and with its separate noiseless extension.
